\documentclass[10pt,a4paper]{amsart}
\usepackage[margin=19mm]{geometry}
\usepackage{amsmath,amssymb,mathtools}
\usepackage[scr=rsfs,cal=euler]{mathalfa}
\usepackage{xfrac}
\usepackage[unicode,hidelinks,bookmarks=false]{hyperref}
\newcommand{\ie}{i.e.\ }
\newcommand{\cf}{cf.\ }
\newcommand{\resp}{respectively\ }
\newcommand{\Subsection}{\subsection}
\providecommand{\nobreakdash}{\nobreak\hbox{-}}
\setlength{\emergencystretch}{2em}
\multlinegap0pt
\usepackage{bm}
\usepackage{bbm}
\usepackage{dsfont}
\usepackage{xspace}
\usepackage{cancel}
\usepackage{mathtools}
\usepackage[shortlabels]{enumitem}
\usepackage{amsthm}
\makeatletter
\@ifundefined{thm}{\newtheorem{thm}{Theorem}}{}
\@ifundefined{prop}{\newtheorem{prop}[thm]{Proposition}}{}
\makeatother
\newcommand{\N}{\mathbb{N}}
\newcommand{\h}{\mbox{${\mathcal{H}}$}}
\newcommand{\leql}{\leq_{\mathcal{L}}}
\newcommand{\sm}{\setminus}
\renewcommand{\j}{\mbox{${\mathcal{J}}$}}
\renewcommand{\l}{\mbox{${\mathcal{L}}$}}
\title[On minimal conditions in semigroups and biacts]{On minimal conditions in semigroups and biacts}
\author{Craig Miller}
\date{Source: arXiv:2506.13617v1 (CC BY 4.0)}
\begin{document}
\maketitle

\noindent\emph{Source and licence.} On minimal conditions in semigroups and biacts. Version arXiv:2506.13617v1 (CC BY 4.0), distributed under the Creative Commons Attribution 4.0 International licence, as recorded in the arXiv licence metadata for this exact version and shown on the abstract page https://arxiv.org/abs/2506.13617v1. Full rights notice: licenses/arxiv-2506.13617-CC-BY-4.0.txt.\par\smallskip

\noindent\textit{Editorial scope.} Counted unit: the Thm whose source label is \texttt{thm:Greenindex,stab} in the article (source0.tex lines 747--754, source label \texttt{thm:Greenindex,stab}), delivered together with the article's own complete original proof of it (source0.tex lines 756--774, one \texttt{proof} environment of 19 lines). No mathematics is written, repaired or re-derived: statement and proof are the article's own text line for line. Required context retained uncounted: prop:subsemigroup,L-classes (lines 723--729). Every definition, display and label of those lines is retained unchanged. Omitted from this package and not redistributed: 1--722, 730--746, 755--755, 775--915. Nothing inside a retained excerpt is omitted or altered: every retained body is delivered exactly as the article's lines are written, so no omission receipt of any kind accompanies this package. Citations: the retained excerpts carry no citation command at all, so no bibliography entry of the article is transcribed and no reference is redistributed. Reference adaptation: none was needed and none was made; every reference inside the delivered unit resolves inside it, so no unresolved reference or citation is produced. Printed slips kept literal: no printed word, symbol, hypothesis or number of the counted statement or of its proof was altered. Layout and class: the article's own class is \texttt{amsart} with packages including fullpage, amsfonts, amssymb, amsmath, mathrsfs, enumitem, tikz, float, subcaption, cite; the wrapper is the lane's pinned \texttt{amsart} editorial wrapper at 10pt on A4, which restates the theorem-like environments the retained text needs and adds the article's own macro definitions that the retained text needs (\texttt{\textbackslash N}, \texttt{\textbackslash h}, \texttt{\textbackslash j}, \texttt{\textbackslash l}, \texttt{\textbackslash leql}, \texttt{\textbackslash sm}), with extra wrapper packages (bm, bbm, dsfont, xspace, cancel, mathtools, enumitem). The delivered numbering is the unit's own. Assets: no retained span contains a figure environment, a graphics-inclusion command, a PDF-page inclusion, a table or a TikZ picture; no image file is copied into this package, the package declares no asset dependency and no image file is redistributed with it. Licence: version arXiv:2506.13617v1 of this article is distributed under the Creative Commons Attribution 4.0 International licence, as recorded in the arXiv licence metadata for this exact version and shown on the abstract page https://arxiv.org/abs/2506.13617v1. A full rights notice is delivered at licenses/arxiv-2506.13617-CC-BY-4.0.txt. Characters: every retained excerpt is pure ASCII, so the delivered engine, XeLaTeX with its default Latin Modern text font, reports no missing character.
\begin{prop}\label{prop:subsemigroup,L-classes}
Let $S$ be a semigroup, and let $T$ be a subsemigroup of $S$ such that ${_T}(S/T)_T$ has finitely many $\l$-classes.
\begin{enumerate}[leftmargin=*]
\item[\emph{(1)}] If $S$ is left stable then $T$ is left stable.
\item[\emph{(2)}] The semigroup $T$ is left stable if and only if ${_T}S_T$ is left stable.
\end{enumerate}
\end{prop}

\begin{thm}\label{thm:Greenindex,stab}
Let $S$ be a semigroup, and let $T$ be a subsemigroup of $S$ of finite Green index.  Then the following are equivalent:
\begin{enumerate}[leftmargin=*]
\item[\emph{(1)}] $S$ is (left) stable;
\item[\emph{(2)}] $T$ is (left) stable;
\item[\emph{(3)}] ${_T}S_T$ is (left) stable.
\end{enumerate}
\end{thm}

\begin{proof}
We prove the result for left stability.  Then, by symmetry, a dual statement holds for right stability, and the result for (two-sided) stability follows.   

Given Proposition \ref{prop:subsemigroup,L-classes}, we only need to prove that (2) implies (1).  So, let $a,b\in S$ with $a\leql b$ and $a\,\j\,b$ in $S.$  Then there exist $s,u,v\in S^1$ such that $a=sb$ and $b=uav.$  Then $a=(su)^iav^i$ and $b=u(su)^iav^{i+1}$ for all $i\in\N.$ 

Suppose first that there are infinitely many $n\in\N$ such that $(su)^n\in S{\sm}T$ or $(su)^na\in S{\sm}T.$  Since $T$ has finite Green index in $S,$ and since $\h\subseteq\l$ and $\l$ is right compatible, it follows that there exist $i,j\in\N$ with $i<j$ such that $(su)^ia\,\l\,(su)^ja$ in ${_T}(S/T)_T$.  Thus, there exists $x\in T^1$ such that $(su)^ia=x(su)^ja.$  It follows that
$$b=u(su)^iav^{i+1}=ux(su)^jav^{i+1}=ux(su)^{j-i-1}(su)^{i+1}av^{i+1}=ux(su)^{j-i-1}a,$$
and hence $a\,\l\,b$ in $S.$

Now suppose that there are only finitely many $n\in\N$ such that $(su)^n\in S{\sm}T$ or $(su)^na\in S{\sm}T.$  Then there exists $m\in\N$ such that $(su)^n,(su)^na\in T$ for all $n\geq m.$  Thus, for each $n\geq m,$ we have $(su)^{m+n}a=(su)^n(su)^ma\in T^1(su)^ma.$  We claim that there exists $N\geq m$ such that $(su)^{m+N}a\,\j\,(su)^ma$ in $T.$  Indeed, for each $n\geq m$ we have 
$$(su)^ma=(su)^m(su)^{m+n}av^{m+n}\in T^1(su)^{m+n}av^{m+n}.$$  
Thus, if $v^{m+n}\in T^1$ for some $n\geq m,$ the claim holds with $N=n.$  Suppose then that $v^{m+n}\in S{\sm}T$ for all $n\geq m.$  Since $T$ has finite Green index in $S,$ it follows that there exist $i,j\geq m$ with $j-i\geq m$ such that $v^i\,\h\,v^j$.  Thus, there exists $y\in T^1$ such that $v^j=v^iy.$  It follows that
$$(su)^ma=(su)^{m+j}av^iy
=(su)^{m+j-i}(su)^iav^iy=(su)^{m+j-i}ay\in(su)^{m+j-i}aT^1,$$
and hence the claim holds with $N=j-i.$  This establishes the claim.  Now, since $T$ is left stable, there exists $t\in T^1$ such that $(su)^ma=t(su)^{m+N}a.$  We then have
$$b=u(su)^mav^{m+1}=ut(su)^{m+N}av^{m+1}
=ut(su)^{N-1}(su)^{m+1}av^{m+1}=ut(su)^{N-1}a\in S^1a.$$
We conclude that $a\,\l\,b$ in $S,$ as required.
\end{proof}

\end{document}
